\chapter{Introduction}
\label{ch:introduction}

Vision-and-Language Navigation (VLN) requires an agent to navigate an unseen 3D environment to reach a target location by following natural language instructions and real-time egocentric visual observations. With their strong vision-language reasoning and generalization capabilities, Multimodal Large Language Models (MLLMs) have emerged as the core architecture for VLN agents, where the ability to infer 3D physical structure from 2D RGB observations is critical to navigation success and efficiency. However, the visual encoders of existing MLLMs suffer from fundamental limitations in 3D spatial reasoning. Feature distillation, which transfers representations from pretrained geometry models into MLLM visual tokens, has gained attention as a promising remedy. Nevertheless, prior studies rely on a single teacher model and a fixed distillation pipeline, leaving the key design conditions for effectively transferring 3D geometric priors to VLN MLLMs systematically unexamined.

To bridge this gap, this study decomposes the impact of geometry-based feature distillation on MLLM navigation performance along three design axes. First, we investigate teacher model selection. Recently proposed feed-forward models and the recurrent model differ markedly in their backbone architectures and 3D decoding mechanisms, raising the question of which geometry model serves as the most effective distillation target for VLN. Second, we examine information flow directionality. Whereas multi-view geometry models attend bidirectionally across all frames to capture cross-view consistency, MLLM decoders enforce a unidirectional information flow via causal attention. By comparing bidirectional and causal variants of the VGGT family and CUT3R, we analyze how this directional mismatch between teacher and student affects distillation performance. Third, we study the temporal stride used to construct the navigation context. Unlike standard 2D encoders that process images independently, multi-view geometry models jointly encode multiple frames; accordingly, we evaluate how changes in inter-frame visual overlap induced by the history sampling interval affect teacher feature quality and downstream VLN performance.

Our systematic analysis yields three key findings. First, under standard final-layer feature distillation, Pi3 which decodes geometry directly from its terminal backbone outputs without multi-layer aggregation achieves substantial performance gains, whereas DPT-head-based models yield only marginal improvements or degrade nDTW on RxR-CE. This indicates that structural alignment between the teacher's decoding head and the distillation target layer is a critical design factor. Second, neither bidirectional nor causal variants exhibit a consistent advantage across the VGGT and CUT3R families, demonstrating that information flow directionality mismatch is not a dominant bottleneck in the current VLN feature distillation pipeline. Third, widening the temporal stride of navigation history improves overall performance by providing longer-term context, yet progressively narrows the additional gains from distillation. This reveals a structural trade-off: while wider strides benefit long-horizon context modeling, they reduce inter-frame visual overlap and weaken the 3D spatial reasoning benefits of multi-view geometry teachers.
